\section{Fundamentos teóricos}

\subsection{Aprendizaje supervisado y generalización}
En clasificación, cada observación contiene predictoras $x_i\in\mathbb{R}^d$ y una
etiqueta discreta $y_i$. El modelo aprende una regla $\hat y=f(x)$ a partir de
datos de entrenamiento, y su utilidad se mide en observaciones que no participaron
en ese aprendizaje \cite{hastie}. Los \emph{hiperparámetros} (por ejemplo, el número
de vecinos o la profundidad de un árbol) controlan el procedimiento de ajuste; los
\emph{parámetros} se estiman con los datos. Más complejidad reduce el error de
entrenamiento, pero puede aumentar la varianza ante muestras nuevas: el modelo
memoriza ruido (sobreajuste). Si el modelo es demasiado simple, no captura la
estructura (subajuste). Por eso hay que separar ajuste, selección y evaluación final
\cite{sklearnleakage,sklearnmetrics}.

\subsection{KNN: vecinos más cercanos}
KNN no construye un modelo explícito: guarda las observaciones de entrenamiento y
clasifica un punto nuevo según las etiquetas de sus $k$ vecinos más próximos
\cite{coverhart,sklearnknn}. Con la distancia de Minkowski,
\[
 d_p(x,z)=\Bigl(\sum_{j=1}^{d}|x_j-z_j|^p\Bigr)^{1/p},
 \qquad
 \hat y=\arg\max_c\sum_{i\in N_k(x)}w_i\,\mathbf{1}(y_i=c),
\]
donde $p=2$ es la distancia euclidiana, $p=1$ la Manhattan, y los pesos $w_i$ son
uniformes o inversamente proporcionales a la distancia. Un $k$ pequeño conserva
detalles locales y es sensible al ruido; uno grande suaviza la frontera y puede
diluir una clase poco frecuente, porque la mayoría de los vecinos pertenecerá a la
clase dominante. Como todas las variables entran en la distancia con el mismo peso,
el escalado es obligatorio y los atributos irrelevantes empeoran el método: con
muchas dimensiones las distancias tienden a parecerse entre sí.

\subsection{Árbol de decisión (DT)}
Un árbol divide recursivamente el espacio de atributos con reglas de la forma
$x_j\le t$. Cada hoja es una región y predice la clase más frecuente de las
observaciones de train que caen en ella. La implementación utilizada,
\texttt{DecisionTreeClassifier}, sigue el algoritmo CART \cite{sklearntree}. Para las
proporciones $p_c$ de cada clase en un nodo, las impurezas habituales son
\[
 G=1-\sum_c p_c^2 \quad\text{(Gini)}, \qquad H=-\sum_c p_c\log p_c \quad\text{(entropía)},
\]
y en cada nodo se elige la división que más reduce la impureza ponderada de los
hijos. Sin restricciones, el árbol crece hasta que cada hoja es pura y memoriza el
entrenamiento. La profundidad máxima y el mínimo de observaciones por hoja limitan
esa complejidad. Sus reglas se pueden leer directamente, pero cambios pequeños en la
muestra producen árboles distintos: es un método de alta varianza.

\subsection{Random Forest (RF)}
Random Forest promedia muchos árboles profundos, cada uno entrenado con una muestra
\emph{bootstrap} de las filas y con un subconjunto aleatorio de atributos candidatos
en cada división \cite{breiman,sklearnensemble}. Cada árbol individual sobreajusta,
pero sus errores están poco correlacionados y el promedio reduce la varianza sin
aumentar mucho el sesgo. \texttt{n\_estimators} fija cuántos árboles se promedian
(más árboles estabilizan la predicción pero no cambian el sesgo), y
\texttt{max\_depth} o \texttt{min\_samples\_leaf} limitan cada árbol. Es habitual que
un bosque alcance exactitud casi perfecta en train sin que eso implique la misma
calidad en datos nuevos.

\subsection{XGBoost}
XGBoost construye un ensamble \emph{aditivo}: cada árbol nuevo corrige los errores
del conjunto anterior siguiendo el gradiente de una función de pérdida, con una
penalización explícita de complejidad \cite{chen,xgboostmodel}:
\[
 F_T(x)=\sum_{t=1}^{T}\eta f_t(x),\qquad
 \mathcal L=\sum_i\ell\bigl(y_i,F_T(x_i)\bigr)+\sum_t\Omega(f_t),
 \qquad \Omega(f)=\gamma L+\tfrac{\lambda}{2}\sum_{j=1}^{L}w_j^2,
\]
donde $L$ es el número de hojas, $w_j$ el valor de cada hoja y $\eta$ la tasa de
aprendizaje. Una $\eta$ menor hace que cada árbol aporte menos y actúa como
regularización; una profundidad menor limita las interacciones que cada árbol
puede representar. Se usa la biblioteca \texttt{xgboost.XGBClassifier}, no el
\texttt{GradientBoostingClassifier} de scikit-learn.

\subsection{AdaBoost}
AdaBoost combina en secuencia aprendices débiles, por defecto árboles de un solo
nivel (\emph{stumps}), y aumenta el peso de las observaciones mal clasificadas para
que el siguiente aprendiz se concentre en ellas \cite{freund,sklearnada}. En la
versión binaria, con error ponderado $\epsilon_t$,
\[
 \alpha_t=\tfrac12\log\frac{1-\epsilon_t}{\epsilon_t},\qquad
 D_{t+1}(i)\propto D_t(i)\,e^{-\alpha_t y_i h_t(x_i)},\qquad
 \hat y=\operatorname{sign}\Bigl(\sum_t\alpha_t h_t(x)\Bigr).
\]
El número de estimadores y la tasa de aprendizaje regulan el ensamble. Si cada stump
apenas supera el azar ($\epsilon_t\approx0.5$), su peso $\alpha_t$ es casi nulo y el
ensamble aprende muy poco; además, concentrar peso en casos difíciles puede
amplificar el ruido de las etiquetas.

\subsection{Evaluación: matriz de confusión y métricas}
La matriz de confusión cuenta $C_{ij}=\#\{y=i,\hat y=j\}$ (filas reales, columnas
predichas). Para cada clase $c$, tratada como positiva frente al resto,
\[
 P_c=\frac{TP_c}{TP_c+FP_c},\qquad
 R_c=\frac{TP_c}{TP_c+FN_c},\qquad
 F1_c=\frac{2P_cR_c}{P_c+R_c},\qquad
 \text{F1 macro}=\frac1{|C|}\sum_c F1_c .
\]
La exactitud $\sum_c C_{cc}/n$ puede ser alta aunque una clase poco frecuente no se
detecte nunca. El F1 macro da el mismo peso a cada clase y por eso es la métrica de
selección de este trabajo: los dos conjuntos están desbalanceados (60/40 y 85/15)
\cite{sklearnmetrics,hegarcia}.

La \textbf{curva ROC} grafica la tasa de verdaderos positivos contra la de falsos
positivos al mover el umbral de decisión sobre la probabilidad estimada. El área
bajo la curva (AUC) es la probabilidad de que el modelo asigne mayor puntuación a un
positivo que a un negativo elegidos al azar: 0.5 equivale al azar y 1 a una
separación perfecta. No depende del umbral 0.5 que se usa para producir la matriz
de confusión \cite{fawcett}.

\subsection{Validación cruzada estratificada}
Una única partición de validación es muy sensible a qué filas caen en ella,
sobre todo con pocos casos de la clase minoritaria. En la validación cruzada de $K$
pliegues, el entrenamiento se divide en $K$ partes; cada configuración se entrena $K$
veces con $K-1$ partes y se evalúa en la restante, y se reporta la media y la
desviación de la métrica \cite{kohavi,sklearncv}. La versión \emph{estratificada}
conserva la proporción de clases en cada pliegue. Todo preprocesamiento (imputación,
codificación, escalado y selección de atributos) se ajusta dentro de cada pliegue,
para que el pliegue de evaluación nunca influya en el modelo. El conjunto de prueba
se separa antes y sólo se usa una vez, con el modelo ya elegido.

\subsection{Selección de atributos por información mutua}
La información mutua $I(X_j;Y)$ mide cuánto se reduce la incertidumbre sobre la clase
al conocer una variable; es cero si son independientes y, a diferencia de la
correlación, también detecta relaciones no lineales \cite{kraskov,ross}. Un método de
filtro como \texttt{SelectKBest} conserva las $k$ variables con mayor $I(X_j;Y)$
\cite{guyon}. Su limitación es que evalúa cada variable por separado: una variable que
sólo es útil en combinación con otra puede recibir una puntuación cercana a cero.

\subsection{Desbalance de clases y ponderación}
Con clases desbalanceadas, un algoritmo que minimiza el error total tiende a favorecer
la clase mayoritaria. Una corrección simple es ponderar cada observación de train con
$w_i=\frac{n}{|C|\,n_{y_i}}$, inversamente proporcional a la frecuencia de su clase,
de modo que ambas clases pesen lo mismo en el ajuste \cite{hegarcia}. Se aplica a DT, RF,
XGBoost y AdaBoost; KNN no admite pesos por observación.

\subsection{Importancia por permutación}
Para interpretar un modelo ya entrenado se permuta al azar una variable en el conjunto
de prueba y se mide cuánto empeora la métrica. Si la caída es grande, el modelo
depende de esa variable; si es nula, la ignora. Al ser independiente del algoritmo,
permite comparar KNN con los métodos basados en árboles \cite{breiman,sklearnperm}.
No demuestra causalidad.

\subsection{Mapa conceptual}
La Figura~\ref{fig:mapa} resume cómo se relacionan los conceptos anteriores dentro del
flujo experimental.

\begin{figure}[H]\centering
\begin{tikzpicture}[
  font=\footnotesize,
  caja/.style={draw, rounded corners=2pt, align=center, inner sep=4pt, anchor=north, fill=#1},
  caja/.default=white,
  flecha/.style={-{Stealth[length=2mm]}, thick, gray!70!black}]
\node[caja=gray!12, text width=13.2cm] (sup) at (0.5,0) {\textbf{Aprendizaje supervisado: clasificación}\\ datos etiquetados $(x_i,y_i)$ $\rightarrow$ regla $\hat y=f(x)$ que generalice a datos nuevos};
\node[caja=blue!8, text width=5.0cm] (inst) at (-3.6,-1.6) {\textbf{Basado en instancias}\\ KNN: vota con los $k$ vecinos;\\ depende de la distancia\\ y del escalado};
\node[caja=green!10, text width=5.6cm] (arb) at (3.7,-1.6) {\textbf{Basado en árboles}\\ DT: reglas $x_j\le t$ elegidas por impureza (Gini, entropía)};
\node[caja=yellow!15, text width=5.0cm] (pre) at (-3.6,-3.75) {\textbf{Preparación (sólo con train)}\\ imputación, one-hot, escalado,\\ selección por información mutua,\\ ponderación de clases};
\node[caja=green!18, text width=3.3cm] (bag) at (1.9,-3.75) {\textbf{Bagging}\\ RF: promedia árboles independientes\\ $\rightarrow$ menos varianza};
\node[caja=orange!15, text width=3.3cm] (boost) at (5.5,-3.75) {\textbf{Boosting}\\ secuencial, corrige errores: AdaBoost (pesos), XGBoost (gradiente)};
\node[caja=purple!10, text width=13.2cm] (val) at (0.5,-6.2) {\textbf{Selección del modelo}: estándar y variantes de un solo cambio $\rightarrow$ validación cruzada estratificada de 5 pliegues $\rightarrow$ ganador por familia con mayor F1 macro};
\node[caja=red!8, text width=13.2cm] (eval) at (0.5,-7.5) {\textbf{Evaluación final en test reservado}: matriz de confusión, precisión, recall y F1 por clase, ROC/AUC\\ \textbf{Interpretación}: importancia por permutación, reglas del árbol, sobreajuste (train vs.\ CV vs.\ test)};
\draw[flecha] (sup.south -| inst.north) -- (inst.north);
\draw[flecha] (sup.south -| arb.north) -- (arb.north);
\draw[flecha] (arb.south -| bag.north) -- (bag.north);
\draw[flecha] (arb.south -| boost.north) -- (boost.north);
\draw[flecha] (inst) -- (pre);
\draw[flecha] (pre.south) -- (pre.south |- val.north);
\draw[flecha] (bag.south) -- (bag.south |- val.north);
\draw[flecha] (boost.south) -- (boost.south |- val.north);
\draw[flecha] (val) -- (eval);
\end{tikzpicture}
\caption{Mapa conceptual: familias de clasificadores, preparación sin fuga de información, selección por validación cruzada y evaluación final.}
\label{fig:mapa}
\end{figure}
